% 04 Apr: AI clean

\section{The Free Church and the Congregation}

Editorial article in “Folkebladet” for December 16, 1896.

Everyone knows that the new movement and *Folkebladet* labored with power and earnestness for the cause of union; and we are now reproached because we are not so zealous for the product of the union, the United Church, as we formerly were for the cause of union.

To this we must answer: just as surely as it is true that the union would have been impossible and inconceivable without the diligent labor of the new movement and of Augsburg for catechism instruction and the congregation, or what is the same thing, for peace and freedom, just so surely is it also true that we have been deceived and disappointed in our expectations. According to the principles laid down in its constitution, the United Church was to have been the body of congregational freedom; in practice it became the exact opposite. The United Church has not once, but time after time, broken its constitution, both in spirit and in letter, precisely at this essential point.

Hence there is now a new struggle all along the line for the congregation and its freedom. And because this cause is altogether firmly and clearly grounded in God’s Word, we are not in the least degree in doubt either about its legitimacy or about its outcome. We know that it will require much suffering; we know also that it may take a long time and pass through many apparent defeats; but for all that it stands no less firm that the cause is good and worth suffering and contending for. The little weak beginning that has already taken place also shows signs of having such vitality in it that it is not altogether strange, after all, if they try to crush it by litigation and other acts of destruction.

And when it is now also subtly hinted that “Augsburg’s Friends” are not even satisfied with the work of “the Conference,” but want something freer and better, then we answer that it was, after all, the intention of the new constitution of the United Church that it should become a freer and better body than the Conference. For that reason, on our side, we had nothing against the fact that the other bodies were unwilling to enter the Conference; but they insisted absolutely that a new body should be formed with a new constitution more in agreement with the principles of the new movement than the constitution of the Conference; and when, then, a new body was formed with a new constitution more in agreement with the principles of the new movement than the constitution of the Conference, and when we were allowed to remain in this new body because we stood firm on the principle of the congregation, then one can hardly expect of us that we should now take a step backward, when all along we have worked for progress in the interest of congregational freedom.

And if we believe that we can learn something little by little and can follow the light of divine guidance, so that we can now see a little more clearly into the true nature of free-church work than before, then it is surely not so great a fault to try to correct one’s errors and to do it better little by little.

If one would then ask how matters stand with the congregation in the Free Church, then our answer is this: exactly and altogether as God’s Word in the New Testament sets forth the congregation and its position. We intend, altogether definitely and quite unsparingly, to conform ourselves to the Word; that is the only authority to which a Lutheran Free Church can bow.

The first thing that then comes forth—as the scriptural principle for the relation between congregation and church body—is this: the congregation is the original, and the church body the derivative; the congregation is the proper, true, apostolic, and free form of God’s kingdom on earth in the time from the outpouring of the Spirit until Christ’s coming; all outward bodies or churches or institutions or associations or societies have, according to God’s Word, only their justification as serving organs for the congregation, not as ruling and commanding authorities.

This principle is the first and most far-reaching of all rules for church polity. The individual local congregation has God’s Word over it, has Jesus Christ as Head and Lord, has God’s Spirit as its Guide; but it has no other authority over it, and it has no right to give away or lend out its Christian freedom to any other. It can, of course, do so, but it is not biblical or Christian to do so; therefore neither is it without responsibility before God to do so. According to God’s Word the congregation is free, but not so free that it can give away that freedom without becoming guilty and suffering harm.

It is altogether clear that the New Testament requires that Christians should not only have an inward fellowship of the Spirit with one another, but that they should also form outward, organized congregations with their officers and their rules. We would depart from the altogether clear requirement of God’s Word if we tried to be so spiritual that we would have no outward congregation. The outward congregation belongs to apostolic and biblical Christianity. And the many who will not belong to any outward congregation because it is not spiritual enough for them—their Christianity suffers from a sorrowful deficiency and frailty, though it is far from us to mean that such frail Christians are not Christians at all.

On the other hand, the same does not apply to what is called a church fellowship. Everyone who knows a little about Christianity knows that the Christian congregations have spiritual fellowship with one another, and we shall later try to elucidate that matter more fully; but God’s Word shows us no outward, organized church body over or under or outside the local congregation that should be any necessity for the believers or for the congregations. According to Scripture we must have outward congregations and congregational officebearers; we are not, on the other hand, required to have church bodies and offices of the church body and officers of the church body. The congregations can form bodies or refrain from forming bodies, and yet be equally Christian.

The very first fundamental principle of the Free Church for church polity is therefore this: that the Christian Church consists of free, independent, and self-governing congregations. The local Christian congregation has an independent authority that applies to itself alone. It owes obedience to no other congregation, nor does it have the right to rule over any other congregation. And what in this matter applies to a single congregation applies to every single one of them and to them all together. Among one hundred congregations one congregation cannot rule; neither can the ninety-nine rule over the one.

So teaches Scripture. We do well to emphasize it sharply and hold it firmly and immovably. Therefore, if there is anyone who is eager to labor in the Free Church, let him begin with the congregation. In this matter it is Alpha and Omega. And since there is so much lacking in our state-church congregations as to what they ought to be, there is at once, and immediately, room for all the work we can do in this one thing: the liberation of the congregation.

Georg Sverdrup
